% Chapter 8, typeset from lectures/L08: the README is the opener and the closing review, and each
% part of Appendix A is a section, in the same order and with the same subsections: A.1 to A.3,
% the machine designed by hand, in design.tex, and A.4 to A.6, the machine in VHDL, in vhdl.tex.
% Appendix B, the exercises and the two course capstones, is the last section.

\chapter{State Machines}
\label{c8:ch}

{\large\itshape Designed by hand first, then expressed in the language.\par}

\begin{chapterbox}{In this chapter}
\begin{itemize}
  \item Finite state machines: states, transitions, inputs and outputs.
  \item Hand-designing a Moore machine: state diagram, state table, Karnaugh-derived logic.
  \item State encoding as a design decision, and why the worked machine uses a Gray code.
  \item The same machine in VHDL: an enumerated type and the \code{case} transition pattern.
  \item Why Moore is the default in this book.
  \item \code{fsm_led} on the DE0-CV: buttons cycling the LED off, blinking, on.
  \item Mealy machines (\secref{c8:a5}), and \code{seq_detect_mealy} to read and run.
  \item The course capstones, as exercises: a serial receiver and a serial transmitter, the two ends
    of one link.
\end{itemize}
\end{chapterbox}

This chapter designs one state machine by hand and then writes another in VHDL. The first is a
Moore machine taken from its specification through a state diagram and a state table to
Karnaugh-derived next-state and output logic, with no step skipped, and then realized as a gate
network in CircuitVerse and stepped through its state cycle by hand: at each step, predict the next
state from the diagram before the clock is advanced. The second is \code{fsm_led}, the three-state
LED controller of \secref{c8:a4}, composed from Chapter~\ref{c4:ch}'s synchronizer and
Chapter~\ref{c7:ch}'s timer. An enumerated \code{state_t} type and a \code{case}-statement process
replace the state table and the Karnaugh maps entirely. It is a different, larger machine than the
hand-designed one, deliberately, so the correspondence you are watching is structural rather than a
transcription: every row you derived by hand in the state table becomes one branch of a
\code{case}, and seeing that correspondence is worth more than either step on its own. The course
then takes \code{fsm_led} onto the board.

Mealy machines are read rather than presented: \secref{c8:a5} covers them, with
\code{seq_detect_mealy} to read and run, and Exercise~\ref{c8:ex:4} has you write one. Reading it
matters more than it might look. Everything else in the chapter is Moore, and Moore is the default
in this book. \Secref{c8:a5} is where you find out what the alternative costs and buys, which is the
only way that default becomes a choice rather than a habit.

The chapter assumes Chapter~\ref{c2:ch}'s Karnaugh maps and Chapter~\ref{c3:ch}'s flip-flops, and
reuses Chapter~\ref{c7:ch}'s \code{timer} unchanged, together with Chapter~\ref{c4:ch}'s
synchronizer widened to two buttons. Of the exercises, the first two are pen-and-paper,
Exercise~\ref{c8:ex:5} is reading and running the worked machine, and the rest are VHDL.
Exercises~\ref{c8:ex:6} and~\ref{c8:ex:7} are the \textbf{course capstones}: the two halves of one
serial link. The receiver composes Chapter~\ref{c4:ch}'s \code{reset_sync},
Chapter~\ref{c5:ch}'s \code{sync} and Chapter~\ref{c7:ch}'s \code{timer} around a state machine
that does its own shifting; the transmitter needs only the first and the last of those, mirrored.
Budget real time for them, and design each state diagram on paper before writing any VHDL. Do the
receiver first: it is the harder of the two, because it has to recover the bit timing it is given
rather than define it. The modules they reuse are worth having open alongside them, and so are
Chapter~\ref{c6:ch}'s \code{serial_rx8} and \code{piso8} for the shift idiom, though neither is
instantiated: the capstones shift inside their own state machines.

% Sections 8.1 to 8.3, from lectures/L08/appendix/a_state_machines.md, parts A.1 to A.3: the
% machine designed by hand.

\section{What Is a Finite State Machine}
\label{c8:a1}
A finite state machine (FSM) is a circuit whose behaviour depends not only on its current inputs but
on a finite amount of \emph{history}, summarized in a single value called its \textbf{state}. At any
moment the machine is in exactly one state; on each clock edge it may move to a new one, decided by
its current state and inputs; and its output is computed from that state.

None of this is new hardware. An FSM combines two things you already know: a \textbf{register}
(Chapter~\ref{c3:ch}) holding the current state, one flip-flop per state bit, and
\textbf{combinational logic} computing the next state and the output from the current state and the
inputs.

What is new is the discipline. Instead of ad hoc logic feeding an ad hoc set of flip-flops, you work
in a fixed order: first draw a \textbf{state diagram}, a few named states with labelled transitions
between them, and only once that is right derive the state table, then the next-state logic, then
the gate network or VHDL.

Every synchronous circuit so far is technically a trivial state machine: a counter's state is its
count, a shift register's is its bit pattern. What sets an FSM apart is that the state is not a
count or a shifted pattern but an arbitrary label (\code{STATE_OFF}, \code{STATE_BLINK}, \dots)
whose meaning you choose and whose next value can depend on the input in any way you like.

Two flavours exist, and the distinction decides what an output may look at:
\begin{itemize}
  \item \textbf{Moore machines}: the output depends \emph{only} on the current state. That is the
    kind designed by hand here.
  \item \textbf{Mealy machines}: the output depends on the current state \emph{and} the current
    input. \Secref{c8:a5} covers these briefly, in VHDL only, so you recognize one when you meet it
    and know the one-cycle difference.
\end{itemize}

\section[Hand-Designing a Moore Machine]{Hand-Designing a Moore Machine: State Diagram, State
Table, and Karnaugh-Derived Logic}
\label{c8:a2}
The same process as every gate network in Chapters~\ref{c1:ch} and~\ref{c2:ch}, applied to a
circuit with memory. The worked example is a small, closed, four-state Moore machine: a circular
counter through a Gray-code sequence, advanced one step per button press.

\subsection{Specification}
\label{c8:spec}
\begin{itemize}
  \item Four states, each represented by two flip-flops \code{\{Q1, Q2\}}:

    \begin{narrowtable}{@{}ll@{}}
      \thead{State} & \thead{\code{\{Q1, Q2\}}} \\
      \midrule
      \code{STATE_0} & \code{00} \\
      \code{STATE_1} & \code{01} \\
      \code{STATE_2} & \code{11} \\
      \code{STATE_3} & \code{10} \\
    \end{narrowtable}
  \item A system clock \code{clock}, and an active-low reset \code{reset_n} returning the machine to
    \code{STATE_0}.
  \item An active-low push button \code{button_n}, from which a synchronized, edge-detected pulse
    \code{X} is derived exactly the way Chapter~\ref{c4:ch}'s double-flop synchronizer does.
    \code{X} is high for exactly one clock cycle on a falling edge of \code{button_n}.
  \item An output \code{Y}, connected to an LED, high only in \code{STATE_3}.
  \item The machine is \textbf{closed}: the state after \code{STATE_3} is \code{STATE_0}.
  \item It is a \textbf{Moore} machine: \code{Y} depends only on \code{\{Q1, Q2\}}, never on
    \code{X}.
\end{itemize}

\noindent The encoding is a Gray code (\code{00 -> 01 -> 11 -> 10 -> 00}), where consecutive states
differ in exactly one bit. That is a deliberate choice, not a coincidence:
\begin{itemize}
  \item With an ordinary binary count (\code{00 -> 01 -> 10 -> 11}), a transition like
    \code{01 -> 10} flips two bits at once, and the two flip-flops will not switch at
    \emph{exactly} the same instant.
  \item That skew is harmless to the machine itself. Both flip-flops update on the same clock edge,
    and nothing samples the state register in between, so the machine never acts on the
    intermediate value. This is the same argument as the critical path of \secref{c5:a3}: as long
    as everything settles before the next edge, what happens during the settling does not matter.
  \item What the skew \emph{can} disturb is anything reading the state combinationally. An output
    decoded straight from the state bits, like \code{Y} below, can glitch briefly while they skew.
    So can a value sampled by a \emph{different} clock domain, which is the case \secref{c4:a3}
    warned about.
  \item A Gray code removes both, since only one bit ever changes per transition. That is why it
    matters most on a counter crossing clock domains, and least on a machine whose output is
    registered.
\end{itemize}
Encoding is a design decision in general, and this is the one place in the book you make it by
hand. Binary uses the fewest flip-flops; Gray minimizes bits changing per transition;
\textbf{one-hot} spends one flip-flop per state so the next-state logic is a single wide OR per
state, which is often fastest on an FPGA where flip-flops are plentiful. Once you write the machine
in VHDL as an enumerated type (\secref{c8:a4}), you hand this decision to the synthesis tool, which
picks an encoding for you and re-encodes freely unless you tell it otherwise.

\subsection{State diagram}
\label{c8:diagram}
\bookimage{0.62}{lectures/L08/appendix/images/fsm_state_diagram.png}{The four-state Gray-code Moore
machine: each state holds on \code{X = 0} and advances on \code{X = 1}, and each bubble reads
\code{Q1Q2/Y}.}{c8:fig:statediagram}

\noindent Each state bubble shows the current state and the output as \code{Q1Q2/Y}, and each arrow
is labelled with the input \code{X} that triggers it. Each state loops back to itself while
\code{X=0}, and a synchronized button pulse (\code{X=1}) advances to the next state, wrapping from
\code{STATE_3} back to \code{STATE_0}.

\needspace{14\baselineskip}
\subsection{State table}
\label{c8:table}
Every combination of the current state and input, giving the next state \code{\{Q1+, Q2+\}} and
output \code{Y}:

\begin{narrowtable}{@{}cccccc@{}}
  \thead{Q1} & \thead{Q2} & \thead{X} & \thead{Q1+} & \thead{Q2+} & \thead{Y} \\
  \midrule
  0 & 0 & 0 & 0 & 0 & 0 \\
  0 & 0 & 1 & 0 & 1 & 0 \\
  0 & 1 & 0 & 0 & 1 & 0 \\
  0 & 1 & 1 & 1 & 1 & 0 \\
  1 & 0 & 0 & 1 & 0 & 1 \\
  1 & 0 & 1 & 0 & 0 & 1 \\
  1 & 1 & 0 & 1 & 1 & 0 \\
  1 & 1 & 1 & 1 & 0 & 0 \\
\end{narrowtable}

\subsection{Karnaugh-derived logic}
\label{c8:kmap}
Treating \code{Q1+} and \code{Q2+} each as a Boolean function of \code{\{Q1, Q2, X\}} and running
them through a Karnaugh map, the same technique as Chapter~\ref{c2:ch}:

\begin{console}
Q1+ = Q1X' + Q2X
Q2+ = Q2X' + Q1'X
\end{console}

\noindent Notice the shape of both:
\begin{itemize}
  \item each is ``hold when \code{X} is \code{0}, take something from the other bit when \code{X} is
    \code{1}.''
  \item the hold halves, \code{Q1X'} and \code{Q2X'}, really are mirror images of each other.
  \item the advance halves are \code{Q2X} and \code{Q1'X}, and that complement on the second one is
    not a typo: it is what makes the sequence a Gray code rather than a plain rotation.
\end{itemize}
The output \code{Y} reduces further, straight from the state table, with no dependency on \code{X}
at all, confirming this is a valid Moore machine:

\begin{console}
Y = Q1Q2'
\end{console}

\subsection{Realizing it as a gate network}
\label{c8:gates}
\code{\{Q1, Q2\}} are each a D flip-flop, exactly like Chapter~\ref{c3:ch}'s register:
\code{Q1+}/\code{Q2+} feed the D inputs and \code{Q1}/\code{Q2} are read back from the Q outputs.
\code{Y} is a single two-input AND gate with one input inverted, reading the current state. That is
the entire machine: two flip-flops and a handful of gates computing \code{Q1+}, \code{Q2+} and
\code{Y} from \code{\{Q1, Q2, X\}}.

\section{From Gate Network to CircuitVerse}
\label{c8:a3}
Before writing any of this in VHDL, realize it by hand in CircuitVerse
(\url{https://circuitverse.org/simulator}), exactly as every gate network in Chapter~\ref{c1:ch}
was: double-flop synchronize \code{button_n} and \code{reset_n}, edge-detect the synchronized button
to produce \code{X}, wire \code{Q1+}/\code{Q2+} into a pair of D flip-flops, and wire \code{Y}
straight into the LED.

\bookimage{1.0}{lectures/L08/appendix/images/fsm_gate_network.png}{The four-state Moore machine as
a gate network in CircuitVerse, with the double-flop synchronization and edge detection of
\code{button_n} on the left and \code{Y} driving the LED on the right.}{c8:fig:gatenetwork}

\noindent Reading it left to right:
\begin{itemize}
  \item \code{button_n} passes through three synchronizer flip-flops
    (\code{button_s1_n}/\allowbreak\code{button_s2_n}/\allowbreak\code{button_s3_n}).
  \item An AND gate derives \code{button_edge_s2} (this circuit's \code{X}) from the second and
    third stages.
  \item The four vertical rails labelled \code{Q1Q2XX'} carry the state bits and the edge pulse
    across to the gates on the right, which implement \secref{c8:a2}'s \code{Q1+}/\code{Q2+}
    equations.
  \item The two labelled flip-flops store \code{Q1}/\code{Q2} and feed the final AND/NOT gates
    producing \code{Y}, wired to \code{led}.
  \item \code{reset_n} is synchronized by its own pair of flip-flops at bottom left, exactly like
    every other synchronized reset in this book.
\end{itemize}
Open the live circuit yourself: \file{lectures/L08/appendix/images/fsm_gate_network.cv}. Step it by
hand, one button press at a time, and confirm it visits \code{STATE_0 -> STATE_1 -> STATE_2 ->
STATE_3 -> STATE_0} in order, with the LED lighting only on \code{STATE_3}.

% Sections 8.4 to 8.6, from lectures/L08/appendix/a_state_machines.md, parts A.4 to A.6: the
% machine in VHDL, Mealy machines, and the board. The two worked examples, fsm_led and
% seq_detect_mealy, are printed in full at the end of the subsections that introduce them.

\section{The Same Machine in VHDL}
\label{c8:a4}
That was the last time in this book you will design a machine by hand. Everything from here writes
it directly in VHDL, which is worth having earned rather than started with: the synthesis tool does
exactly what you just did on paper, and when a state machine misbehaves you will be debugging the
diagram, not the VHDL.

The translation replaces the explicit state table and Karnaugh equations with an \textbf{enumerated
type} naming each state, a \textbf{case statement} describing each state's transitions, and the
synthesis tool left to work out the flip-flop encoding and gate-level logic.

\subsection{The enumerated state type}
\label{c8:enum}
Suppose a machine with three states controlling an LED: \code{STATE_OFF}, \code{STATE_BLINK}
(blinking at a fixed rate) and \code{STATE_ON}. VHDL declares exactly those three values as a new
type:

\begin{vhdlcode}
type state_t is (STATE_OFF, STATE_BLINK, STATE_ON);
signal state: state_t;
...
state <= STATE_ON;
\end{vhdlcode}

\subsection{The case-statement pattern}
\label{c8:case}
The machine is almost always a synchronous process updating \code{state} on the rising edge via a
\code{case}, one branch per state, with an asynchronous reset returning it to a known starting
state. Assuming \code{to_next_state} and \code{to_prev_state} are synchronized, edge-detected pulses
produced exactly like \code{X} in \secref{c8:a2} and \secref{c8:a3}:

\begin{vhdlcode}
process(clock, reset_s2_n) is
begin
    if (reset_s2_n = '0') then
        -- On reset, go to STATE_OFF.
        state <= STATE_OFF;
    elsif (rising_edge(clock)) then
        case (state) is
            when STATE_OFF =>
                if (to_next_state = '1') then
                    state <= STATE_BLINK;
                elsif (to_prev_state = '1') then
                    state <= STATE_ON;
                end if;
            -- STATE_BLINK and STATE_ON follow the same shape.
            when others =>
                state <= STATE_OFF;
        end case;
    end if;
end process;
\end{vhdlcode}

\noindent\code{case (state) is ... end case;} must cover every value \code{state_t} can take, and
\code{when others} keeps it exhaustive if you add a state later and forget a branch.

What \code{when others} is \emph{not} is a safety net against a corrupted state register. Once
every state has its own branch, as in \file{fsm_led.vhd} below (the listing above abbreviates two
of them away, so \emph{its} \code{when others} is reachable and would wrongly reset the machine from
\code{STATE_BLINK} or \code{STATE_ON}), every value of \code{state_t} is already named above it, so
the branch is unreachable in the source and synthesis normally optimizes it away rather than
building recovery logic. Once Quartus re-encodes the machine, most bit patterns are illegal, and
recovering from one requires Quartus's Safe State Machine setting rather than anything you can
write in VHDL.

\subsection{Worked example: \code{fsm_led}}
\label{c8:fsmled}
\file{lectures/L08/fsm_led/fsm_led.vhd} implements exactly that three-state machine, extended to two
buttons: \code{button_n(0)} advances, \code{button_n(1)} goes back, the states are closed in a loop,
and \code{STATE_BLINK} toggles the LED every 100\,ms.

Three subcomponents are reused, all of them modules you wrote yourself:
\begin{itemize}
  \item \code{reset_sync}: Chapter~\ref{c4:ch}'s reset synchronizer, unchanged
    (Exercise~\ref{c4:ex:8}).
  \item \code{button_sync}: Chapter~\ref{c4:ch}'s button synchronizer, also unchanged, but
    instantiated with \code{generic map(2)} so it handles both buttons at once, producing a two-bit
    pulse vector. This is the payoff of the generic: the file is identical to the one
    Chapters~\ref{c4:ch} and~\ref{c7:ch} instantiate with \code{1} (Exercise~\ref{c4:ex:9}).
  \item \code{timer}: Chapter~\ref{c7:ch}'s timer, unchanged (Exercise~\ref{c7:ex:4}).
    \code{fsm_led} declares its own \code{TIMER_TICK_COUNT} generic and passes it straight through,
    defaulting to \code{5_000_000} for 100\,ms at \code{50 MHz}, and only enables the timer while in
    \code{STATE_BLINK}.
\end{itemize}
None of the three is shipped in \file{fsm_led/}: copy your own in before building the example. By
this point in the book that is four chapters' worth of your own modules composed into one design,
which is the thing worth noticing about it.

Two synchronous processes implement the machine:
\begin{itemize}
  \item \code{STATE_PROCESS} is the \code{case} transition logic above, driven by
    \code{to_next_state}/\allowbreak\code{to_prev_state}, derived combinationally from
    \code{button_edge_s2}.
    Each requires the \emph{other} button's pulse to be \code{'0'} in the same cycle, so pressing
    both close enough together that their pulses land on one clock cycle moves the machine nowhere
    rather than picking a winner arbitrarily. An ambiguous input deserves a defined answer, and
    ``do nothing'' is the one that cannot surprise you.
  \item \code{LED_PROCESS} drives \code{led_s} purely from the current state: off in
    \code{STATE_OFF}, on in \code{STATE_ON}, toggling on every timer timeout in \code{STATE_BLINK}.
    Because it never reads \code{button_n} or \code{button_edge_s2} directly, only \code{state},
    this is a genuine Moore machine: the LED reacts to a press only indirectly, through the state it
    caused.
\end{itemize}
This is the whole of \file{lectures/L08/fsm_led/fsm_led.vhd}:

\vhdlfile{lectures/L08/fsm_led/fsm_led.vhd}

\section{Mealy Machines: The One-Cycle Difference}
\label{c8:a5}
Everything so far has been a \textbf{Moore} machine, deliberately: Moore is the safer default and
the form you design by hand. A \textbf{Mealy} machine relaxes the rule that defines a Moore machine:
the output may depend on the current state \emph{and} the current input, computed combinationally.
It is not stored in a register, and it changes within the same clock cycle the input does.

\subsection{The behaviour: detect two consecutive \code{1} bits}
\label{c8:mealydetect}
A serial input \code{din} carries one bit per clock cycle, and \code{y} should indicate that the
last two bits were both \code{1}, with overlapping matches counting, so \code{111} reports a match
twice.

As a Mealy machine this needs only two states, because the output can react to the input
\emph{before} it is reflected in the next state: \code{STATE_IDLE} (the previous bit was \code{0},
or this is the first) and \code{STATE_ONE} (the previous bit was \code{1}).

\begin{narrowtable}{@{}lcll@{}}
  \thead{State} & \thead{\code{din}} & \thead{Next state} & \thead{\code{y}} \\
  \midrule
  \code{STATE_IDLE} & 0 & \code{STATE_IDLE} & 0 \\
  \code{STATE_IDLE} & 1 & \code{STATE_ONE} & 0 \\
  \code{STATE_ONE} & 0 & \code{STATE_IDLE} & 0 \\
  \code{STATE_ONE} & 1 & \code{STATE_ONE} & \textbf{1} \\
\end{narrowtable}

\noindent The state process is \secref{c8:a4}'s \code{case} pattern; only the output differs. This
is \file{lectures/L08/seq_detect_mealy/seq_detect_mealy.vhd}, and its output assignment is a single
concurrent statement, entirely outside any clocked process:

\begin{vhdlcode}
y <= '1' when (state = STATE_ONE and din = '1') else '0';
\end{vhdlcode}

\noindent That is a \textbf{conditional signal assignment}, and it is the last piece of VHDL syntax
this book introduces. It is a concurrent statement, so it lives in the architecture body rather than
inside a process, and it describes a multiplexer: pick the first value whose condition holds,
otherwise the value after \code{else}. It is the same hardware you would get from an
\code{if}/\code{else} inside a combinational process, written in one line instead of five, and it is
worth having because a one-line output assignment makes it obvious at a glance what the output
depends on.

\begin{vhdlcode}
-- These two describe exactly the same multiplexer.
x <= a when sel = '1' else b;

process(sel, a, b) is
begin
    if (sel = '1') then
        x <= a;
    else
        x <= b;
    end if;
end process;
\end{vhdlcode}

\noindent Always give it an \code{else}. Leave one out and you have described a signal that keeps
its old value when no condition holds, which is a latch, exactly as in \secref{c5:a3}.
\code{fsm_led} uses this form three times, and Exercise~\ref{c8:ex:4} asks you to use it.

\textbf{That one line is the whole distinction.} It reads \code{din}, so it is Mealy. Move the same
condition inside the clocked process, or drop \code{din} from it, and you have a Moore machine
instead.

This is the whole of \file{lectures/L08/seq_detect_mealy/seq_detect_mealy.vhd}:

\vhdlfile{lectures/L08/seq_detect_mealy/seq_detect_mealy.vhd}

\subsection{What it costs, and what it buys}
\label{c8:mealycost}
The Moore version of this detector needs a \textbf{third} state, \code{M_TWO}, existing purely to
remember for one cycle that the match happened, so the output has something to read that does not
depend on the input. Feed both machines \code{1, 1, 0, 1, 1, 1}:

\begin{narrowtable}{@{}cclclc@{}}
  \thead{Cycle} & \thead{\code{din}} & \thead{Mealy state (before)} & \thead{Mealy \code{y}} &
    \thead{Moore state (before)} & \thead{Moore \code{y}} \\
  \midrule
  1 & 1 & \code{STATE_IDLE} & 0 & \code{M_IDLE} & 0 \\
  2 & 1 & \code{STATE_ONE} & \textbf{1} & \code{M_ONE} & 0 \\
  3 & 0 & \code{STATE_ONE} & 0 & \code{M_TWO} & \textbf{1} \\
  4 & 1 & \code{STATE_IDLE} & 0 & \code{M_IDLE} & 0 \\
  5 & 1 & \code{STATE_ONE} & \textbf{1} & \code{M_ONE} & 0 \\
  6 & 1 & \code{STATE_ONE} & \textbf{1} & \code{M_TWO} & \textbf{1} \\
\end{narrowtable}

\noindent The Moore \code{y} column is exactly the Mealy \code{y} column shifted down one row. That
is the entire trade: Mealy reports the match the \emph{same} cycle the second \code{1} arrives where
Moore takes one more, and Mealy needs 2 states where Moore needs 3.

\subsection{Which to reach for}
\label{c8:whichone}
Moore, almost always, and it is what the rest of this book uses:
\begin{itemize}
  \item its output depends on no primary input, so it cannot follow a noisy one.
  \item a Mealy output can in principle glitch combinationally if its inputs aren't already clean.
  \item note that ``Moore'' on its own does not mean glitch-free. An output decoded combinationally
    from several state bits, like \secref{c8:a2}'s \code{Y = Q1Q2'}, can still glitch while those
    bits skew on a transition. What makes an output glitch-free is \emph{registering} it, the way
    \code{fsm_led}'s \code{LED_PROCESS} does, and that is what you want on anything driving a wire
    off the chip. A serial \code{tx} line decoded straight from the state bits will emit spurious
    edges.
  \item a machine whose output depends on fewer things is a machine with fewer ways to be wrong.
\end{itemize}
Reach for Mealy only when you specifically need the lower latency or the smaller state count, and
know that you are trading a synchronous output for a combinational one when you do.

\section{FPGA Demonstration}
\label{c8:a6}
\code{fsm_led} synthesizes and runs exactly the way every VHDL design has since
Chapter~\ref{c1:ch} (Appendix~\ref{app:quartus}): create a Quartus Prime Lite project targeting the
DE0-CV's \code{5CEBA4F23C7N}, add the \file{.vhd} files for the example and its subcomponents,
compile, assign each port to a physical pin in the Pin Planner, then recompile and program the
board.

For \code{fsm_led}, assign \code{clock} to the 50\,MHz oscillator pin, \code{reset_n} and
\code{button_n(1 downto 0)} to three push buttons, and \code{led} to an onboard LED. Then confirm
that pressing ``next'' repeatedly cycles the LED \code{off -> blinking -> on -> off}, that
``previous'' runs the other way, and that reset returns it to off.

\code{seq_detect_mealy} is deliberately \textbf{not} demonstrated on the board. It consumes one bit
per clock cycle, and a switch flipped by hand holds its value for millions of cycles, so the only
thing an LED can show is the degenerate case: \code{din} held high, \code{y} high from the second
cycle onward. The one-cycle distinction the whole of \secref{c8:a5} is about happens far too fast
to see. Run its testbench instead, which is the general lesson rather than an exception:

\begin{shell}
cd lectures/L08/seq_detect_mealy
cp ../../L04/exercises/reset_sync/reset_sync.vhd .     # the one you wrote in L04
ghdl -a --std=93 reset_sync.vhd seq_detect_mealy.vhd seq_detect_mealy_tb.vhd
ghdl -e --std=93 seq_detect_mealy_tb
ghdl -r --std=93 seq_detect_mealy_tb --assert-level=error --stop-time=10ms
\end{shell}

\noindent Exact pin numbers depend on which DE0-CV buttons and LEDs you choose, assigned the same way
as for every design since Chapter~\ref{c1:ch}. Refer to the DE0-CV's own documentation for its
physical pinout if you need to confirm which header pin corresponds to which labelled button, switch
or LED.

% Section 8.7, from lectures/L08/README.md: the objectives, the questions to test yourself, and the
% README's "Next lecture", which in the last lecture is the end of the course.

\needspace{10\baselineskip}
\section{Review}
\label{c8:sec:review}

\subsection*{What you should be able to do now}
After this chapter, you should be able to:
\begin{itemize}
  \item Explain what a finite state machine is, and identify its states, transitions, inputs and
    outputs in a given circuit.
  \item Design a Moore machine by hand, from specification through state diagram and state table to
    Karnaugh-derived equations, and build the result in CircuitVerse.
  \item Model the same machine in VHDL with an enumerated \code{state_t} type and a \code{case}
    process.
  \item Recognize a Mealy machine on sight and say what its output costs and buys, one clock cycle
    of latency against one extra state, and why Moore is the default in this book.
  \item Compose a state machine with previously built subcomponents rather than rewriting them. By
    the capstone that means three earlier chapters' modules around a machine that is the only new
    logic, drawing on the shift and counter idioms of one more.
\end{itemize}

\needspace{6\baselineskip}
\subsection*{Questions to test yourself}
You should be able to explain:
\begin{enumerate}
  \item What makes a circuit a state machine, rather than just a register with some logic around
    it?
  \item In the hand-designed machine, what exactly is \code{X}, and why must it be a one-cycle
    pulse rather than the button's level?
  \item Why is that machine's state encoding a Gray code? What could go wrong with an ordinary
    binary count?
  \item A counter and a shift register are technically state machines too. What distinguishes them
    from the machine designed here?
  \item Which signals is a Moore output allowed to depend on, and a Mealy output? Given a single
    line of VHDL driving an output, how do you tell which you are looking at?
  \item In \secref{c8:a5}'s comparison, why does the Moore version need an extra state, and why does
    its output lag the Mealy version's by exactly one clock cycle?
  \item In \file{fsm_led.vhd}, why does \code{LED_PROCESS} only ever read \code{state}, never
    \code{button_edge_s2} or \code{button_n} directly? What would change if it did?
  \item Why is \code{din} in \code{seq_detect_mealy} not passed through a synchronizer, unlike
    \code{reset_n}?
  \item What does \code{when others} buy you in a state-transition \code{case}, given that
    \code{state_t} has only the values you declared?
  \item In the capstone receiver, why does the timer run at sixteen times the baud rate rather than
    once per bit, and why does sampling in the middle of a bit make the design tolerate a
    transmitter whose clock is slightly different from yours?
  \item The capstone transmitter uses a whole-bit timer and needs no mid-bit anything. What does it
    have that the receiver doesn't, that makes the difference?
\end{enumerate}

\needspace{6\baselineskip}
\subsection*{Where this leaves you}
This is the final chapter of \emph{Digital Design with VHDL}. Over eight chapters the book moved
from gates and Karnaugh maps, through registers, metastability, and the VHDL language, to counters,
shift registers and timers, and now finite state machines, at every step reasoning about a circuit
by hand before writing it in synthesizable VHDL and watching it run on a real FPGA. Every pattern
built along the way, the double-flop synchronizer, the \code{process}/\code{case} coding style, and
now the Moore state machine, was reused again in this chapter's own examples, and carries forward
unchanged into whatever you build next.

% Section 8.8, from lectures/L08/appendix/b_exercises.md, ending with the two course capstones,
% whose ### headings are \task{...} inside Exercises 8.6 and 8.7.

\section{Exercises}
\label{c8:sec:exercises}

\begin{aside}
\textbf{How to work these.} Exercises~\ref{c8:ex:1} and~\ref{c8:ex:2} are pen-and-paper: they are
about getting a state machine right \emph{before} it becomes code, and nothing in them needs VHDL,
GHDL, or a board. Each exercise is tagged with the kinds of work it asks for: \kinds{Paper} by hand,
\kinds{CircuitVerse} in the browser, \kinds{VHDL} in a module of your own.

\textbf{How to check your work.} From Exercise~\ref{c8:ex:4} on, every exercise that asks you to
write a VHDL module ships a self-checking testbench under \file{lectures/L08/exercises} in the
companion repository, and nothing else; Exercise~\ref{c8:ex:3} modifies a module you already have
and is checked by hand, as it says. Assembling the rest of the directory is part of the work: the
modules you wrote in earlier chapters (\code{reset_sync}, \code{button_sync}, \code{sync},
\code{timer}) are copied in from wherever you wrote them. Each exercise prints the exact commands it
needs. Write your module in its \file{exercises/<module>} directory, using the entity name and the
\textbf{port order} the exercise specifies, then run it with GHDL; \secref{c2:app:c} has the three
commands.

Exercises~\ref{c8:ex:6} and~\ref{c8:ex:7} are the course capstones: the two ends of one serial
link, a receiver and a transmitter. Each composes modules from four different chapters, so each
needs more than one \code{ghdl -a} argument, and so do Exercises~\ref{c8:ex:4} and~\ref{c8:ex:5}.
Where that is the case, the exercise prints the command it needs.

No FPGA board is needed for any exercise. The Quartus synthesis and board-programming steps are
demonstrated in the course, and shown in Appendix~\ref{app:quartus}; your job is to get the VHDL
right, and the testbench is how you confirm it. Short answers to the parts you work on paper are in
Appendix~\ref{app:answers}; reference solutions to the VHDL are published in the companion
repository, under \file{lectures/L08/solutions}. Try each exercise yourself before looking at
either.
\end{aside}

% ------------------------------------------------------------------------------------------
\exerciseset{State Diagrams and State Tables}
\label{c8:set:1}

\exercise{A binary-count Moore machine}{Paper, CircuitVerse}
\label{c8:ex:1}
Design a new four-state Moore machine, \code{\{Q1, Q2\}}, controlling an LED via a single
synchronized, edge-detected button pulse \code{X} (produced exactly like \code{X} in the worked
example of \secref{c8:a2}).

It differs from that worked example in two ways:
\begin{itemize}
  \item The state order is the ordinary binary count \code{00 -> 01 -> 10 -> 11 -> 00},
    \textbf{not} the Gray code used in \secref{c8:a2}.
  \item The LED (\code{Y}) lights in \textbf{two} states instead of one:
    \begin{itemize}
      \item Whenever \code{Q1 = 1} (i.e.\ in both \code{10} and \code{11}).
    \end{itemize}
\end{itemize}

\begin{parts}
  \item Draw the state diagram:
    \begin{itemize}
      \item Four states.
      \item One transition arrow per state for \code{X = 1}.
      \item One self-loop per state for \code{X = 0}.
    \end{itemize}
  \item Fill in the full state table (all eight rows of \code{\{Q1, Q2, X\}}).
  \item Derive \code{Q1+}, \code{Q2+}, and \code{Y} via Karnaugh maps, the same way \secref{c8:a2}
    derives them for the Gray-code version.
  \item Realize the gate network by hand:
    \begin{itemize}
      \item Two D flip-flops for \code{\{Q1, Q2\}}.
      \item Gates computing \code{Q1+}, \code{Q2+}, and \code{Y}.
      \item Build and simulate it in CircuitVerse, double-flop synchronizing and edge-detecting the
        button input exactly as in \secref{c8:a3}.
    \end{itemize}
  \item Verify it in CircuitVerse, against your own state table. There is no testbench for this
    machine, and there is not meant to be: the state table you derived in \textbf{b)} \emph{is} the
    specification, and checking a circuit against a specification you worked out yourself is the
    whole point of the exercise.

    Step the clock by hand and check, in this order:
    \begin{itemize}
      \item \textbf{The advance path.} Pulse \code{X} once per step and walk a full lap,
        \code{00 -> 01 -> 10 -> 11 -> 00}. Before each edge, read the expected next state off your
        table. After the edge, read \code{\{Q1, Q2\}} off the flip-flop outputs. All four
        transitions must agree.
      \item \textbf{The hold path.} In each of the four states, step the clock with \code{X = 0}
        and confirm the machine stays put. That is four more checks, and it is the half people
        skip; a machine that advances correctly but does not hold is a machine whose \code{X'}
        terms are wrong.
      \item \textbf{The output.} In every state, check \code{Y} against your table: high in
        \code{10} and \code{11}, low in \code{00} and \code{01}. Do this as its own pass, not
        while you are watching the state bits.
    \end{itemize}
    Where they disagree tells you which gate to look at. A wrong \code{Q1+} shows up only on the
    transitions where \code{Q1} was supposed to change. A wrong \code{Y} shows up with the state
    sequence still perfectly correct, which is why it is worth a separate pass.
\end{parts}
\note[Tip] Since the state order changed from Gray code to binary, expect \textbf{two} transitions
where both state bits flip at once: \code{01 -> 10} and \code{11 -> 00}. Per \secref{c8:a2} that is
harmless to the machine itself, since nothing samples the state register between edges, and
CircuitVerse's ideal simulation will not show you anything either way.

It is harmless to \code{Y} here too, and for a reason worth naming: \code{Y} is high exactly when
\code{Q1 = 1}, so it is one flip-flop's output taken straight to the LED, with no gate between them
to produce a glitch while the two bits skew. Contrast \secref{c8:a2}'s machine, whose
\code{Y = Q1 * Q2'} \textbf{is} decoded from both bits and can therefore glitch briefly on real
hardware. Ask yourself what would change here if the LED lit in only one state, say \code{10}, and
what you would do about it if \code{Y} were driving something that cared.

% ------------------------------------------------------------------------------------------
\exerciseset{Sequence Detection}
\label{c8:set:2}

\exercise{A Moore \code{1, 0, 1} detector}{Paper}
\label{c8:ex:2}
Design a \textbf{Moore} machine that detects the sequence \code{1, 0, 1} arriving one bit per clock
cycle on an input \code{din}, raising its output \code{y} once the third bit of a match has been
seen.
\begin{itemize}
  \item Design it \textbf{non-overlapping}: after a full match, start looking for a fresh
    \code{1, 0, 1} from scratch, rather than reusing the trailing bits. Exercise~\ref{c8:ex:4} turns
    this machine into VHDL and its shipped testbench checks for it, so an overlapping design fails
    there.
\end{itemize}

\begin{parts}
  \item Draw the state diagram:
    \begin{itemize}
      \item One state per amount of the pattern matched so far.
      \item An arrow out of every state for \code{din = 0} \emph{and} for \code{din = 1}; a sequence
        detector is never idle, it is always either advancing or falling back.
      \item Mark which state drives \code{y = 1}.
    \end{itemize}
  \item Fill in the state table, one row per (state, \code{din}) combination.
  \item Trace it. Pick a six-bit sequence for \code{din} containing at least one match, and record,
    cycle by cycle:
    \begin{itemize}
      \item The state before the edge.
      \item The value of \code{din}.
      \item The output \code{y}.
    \end{itemize}
  \item The trap in this machine is the fall-back arrows, not the forward ones. For each state, say
    what happens on the input that does \emph{not} advance the match:
    \begin{itemize}
      \item Which of them return to the start?
      \item Is there any input that should send you back to a \emph{partially} matched state rather
        than all the way to the beginning? Explain your answer for this pattern.
    \end{itemize}
\end{parts}
\note[Tip] count your states before you draw them. ``Nothing matched yet'', ``saw \code{1}'',
``saw \code{10}'', ``saw \code{101}'' is the natural decomposition, and the last of those is the
one that drives \code{y}.

\textbf{Keep your answer.} Exercise~\ref{c8:ex:4} below implements this same detector in VHDL, as a
Mealy machine, which, per \secref{c8:a5}, needs one state fewer than what you just drew. Working out
which of your states disappears, and why, is part of that exercise.

% ------------------------------------------------------------------------------------------
\exerciseset{Finite State Machines in VHDL}
\label{c8:set:3}

\exercise{A fast-blink state}{VHDL}
\label{c8:ex:3}
Extend \code{fsm_led} with a fourth state, \code{STATE_FAST_BLINK}, blinking at 20\,ms instead of
\code{STATE_BLINK}'s 100\,ms. Follow the worked example in \secref{c8:a4}.

Work on a \textbf{copy} of the module, not on \file{lectures/L08/fsm_led/fsm_led.vhd} itself. Copy
the whole directory somewhere of your own first, along with the three modules of your own that it
instantiates:

\begin{shell}
cp -r lectures/L08/fsm_led /tmp/fsm_led_fast
cp lectures/L04/exercises/reset_sync/reset_sync.vhd /tmp/fsm_led_fast     # the three you wrote
cp lectures/L04/exercises/button_sync/button_sync.vhd /tmp/fsm_led_fast   # yourself
cp lectures/L07/exercises/timer/timer.vhd /tmp/fsm_led_fast
cd /tmp/fsm_led_fast
\end{shell}

\noindent There is no testbench to run here, but the four-state machine should still analyze, and
that is worth doing after every edit: it catches a missing \code{case} branch, or a \code{state_t}
value you added in one process and not the other.

\begin{shell}
ghdl -a --std=93 reset_sync.vhd button_sync.vhd timer.vhd fsm_led.vhd
\end{shell}

\noindent The reason is worth knowing, because it is a real hazard and not busywork. \code{fsm_led}
ships a reference testbench, \file{fsm_led_tb.vhd}, and Exercise~\ref{c8:ex:5} asks you to run it.
That testbench asserts the \textbf{three}-state cycle: it presses the button once from
\code{STATE_BLINK} and checks that the LED goes solid, because in the machine as written the next
state is \code{STATE_ON}. In your four-state version that press lands in \code{STATE_FAST_BLINK},
the LED is still blinking, and the assertion fails. Nothing about the failure tells you which of the
two exercises it belongs to. Keep the reference machine intact and that question never comes up.

There is no testbench for the four-state machine. Check it the way \secref{c8:a2} had you check a
state diagram: predict the sequence of states for a run of button presses, then read your
\code{case} branches back and confirm they agree.

You'll need to:
\begin{itemize}
  \item Add the new value to \code{state_t}.
  \item Add a branch to each \code{case} statement, in both \code{STATE_PROCESS} and
    \code{LED_PROCESS}.
  \item Handle the two blink rates, by either:
    \begin{itemize}
      \item Adding a second timer instance, or
      \item Making the existing timer's tick count switch based on the current state.
    \end{itemize}
\end{itemize}

\exercise{A Mealy \code{1, 0, 1} detector}{VHDL}
\label{c8:ex:4}
Implement the \code{1, 0, 1} detector from Exercise~\ref{c8:ex:2} in VHDL, in a new directory
\file{seq_detect_101_mealy}, this time as a \textbf{Mealy} machine.

Start from your Moore state diagram and convert it:
\begin{itemize}
  \item Your final ``saw \code{101}'' state exists only to hold the output high for one cycle.
  \item A Mealy machine doesn't need it: the output condition can read \code{din} directly in the
    state that has already matched \code{1, 0}.
  \item Drop that state, and move its output into the condition on \code{y}.
\end{itemize}
Then follow the \code{type state_t is (...)} / \code{case} pattern from \secref{c8:a4} and
\secref{c8:a5}:
\begin{itemize}
  \item Reuse \code{reset_sync} for \code{reset_n}, exactly as \file{seq_detect_mealy.vhd} does.
    There is no button in this design, so \code{button_sync} is simply not instantiated.
  \item Make sure the output assignment for \code{y} is a single concurrent statement, outside any
    clocked process, depending on both \code{state} and \code{din}, not just \code{state}.
    \begin{itemize}
      \item Otherwise you've quietly built a Moore machine instead, and its output will arrive one
        cycle late, which is exactly what the testbench checks.
    \end{itemize}
\end{itemize}

\modulefigure{seq_detect_101_mealy}

\selfcheck name your entity \code{seq_detect_101_mealy}, with ports \code{clock}, \code{reset_n},
\code{din} (in) and \code{y} (out), declared in that order to match \code{seq_detect_mealy}. A
testbench is in \file{exercises/seq_detect_101_mealy}; it checks a non-overlapping ``101'' detector.
Your module instantiates \code{reset_sync}, which is not shipped there because you wrote it in
Chapter~\ref{c4:ch}'s Exercise~\ref{c4:ex:8}: copy your own in. This is the first exercise that
needs more than one file on the \code{ghdl -a} line, with the subcomponent first so GHDL sees it
before the module that instantiates it:

\begin{shell}
cd lectures/L08/exercises/seq_detect_101_mealy
cp ../../../L04/exercises/reset_sync/reset_sync.vhd .     # the one you wrote in L04
ghdl -a --std=93 reset_sync.vhd seq_detect_101_mealy.vhd seq_detect_101_mealy_tb.vhd
ghdl -e --std=93 seq_detect_101_mealy_tb
ghdl -r --std=93 seq_detect_101_mealy_tb --assert-level=error --stop-time=10ms
\end{shell}

% ------------------------------------------------------------------------------------------
\exerciseset{Reading the Worked Machine}
\label{c8:set:4}

\exercise{Accounting for \code{fsm_led}}{Paper}
\label{c8:ex:5}
\code{fsm_led} is demonstrated on the DE0-CV in the course (see \secref{c8:a6}). This exercise
checks that you can account for what the demonstration shows, and it ships a reference testbench
you can run yourself. It reuses your two Chapter~\ref{c4:ch} synchronizers and your
Chapter~\ref{c7:ch} \code{timer}, none of which it ships, so copy them in first:

\begin{shell}
cd lectures/L08/fsm_led
cp ../../L04/exercises/reset_sync/reset_sync.vhd .     # the three you wrote yourself
cp ../../L04/exercises/button_sync/button_sync.vhd .
cp ../../L07/exercises/timer/timer.vhd .
ghdl -a --std=93 reset_sync.vhd button_sync.vhd timer.vhd \
                 fsm_led.vhd fsm_led_tb.vhd
ghdl -e --std=93 fsm_led_tb
ghdl -r --std=93 fsm_led_tb --assert-level=error --stop-time=10ms
\end{shell}

\begin{parts}
  \item Work out the blink rate from the source, not from the demonstration:
    \begin{itemize}
      \item \code{fsm_led} instantiates \code{timer} with \code{TIMER_TICK_COUNT} defaulting to
        \code{5_000_000}, on a \code{50 MHz} clock.
      \item How long is one timeout period?
      \item \code{STATE_BLINK} toggles the LED on each timeout. What is the resulting
        \textbf{blink} frequency, and why is it half what you might first write down?
    \end{itemize}
  \item The testbench overrides \code{TIMER_TICK_COUNT} with a much smaller value than the default:
    \begin{itemize}
      \item Find that override in \file{lectures/L08/fsm_led/fsm_led_tb.vhd}.
      \item Explain why a testbench cannot reasonably use the real value, and roughly how long a
        simulation of one real blink period would have to run.
    \end{itemize}
  \item Account for the state machine's behaviour:
    \begin{itemize}
      \item All three states are reachable in both directions. Trace the presses of
        \code{button_n(0)} and \code{button_n(1)} that visit every state and return to
        \code{STATE_OFF}.
      \item Press both buttons on the same clock cycle. What does \code{fsm_led} do, and which two
        lines of the architecture decide that?
      \item \code{timer_enable} is driven from \code{state}, so the timer only runs in
        \code{STATE_BLINK}. Note that the \code{timer} \textbf{pauses} rather than clears when
        disabled (\secref{c7:a3}). What does that mean for the moment you re-enter
        \code{STATE_BLINK} after leaving it, and would holding \code{timer_enable} high permanently
        be better or worse?
    \end{itemize}
\end{parts}

% ------------------------------------------------------------------------------------------
\exerciseset{Capstone: A Serial Receiver}
\label{c8:set:5}

\exercise{The serial receiver}{Paper, VHDL}
\label{c8:ex:6}
This is the first half of the course capstone, and it composes four chapters into one design.
Exercise~\ref{c8:ex:7} builds the other end of the same serial link, so do not stop here. Nothing in
it is a new idea; all of it is wiring together modules you already have, around a state machine you
design yourself:
\begin{itemize}
  \item Chapter~\ref{c4:ch}'s \code{reset_sync}, synchronizing the reset.
  \item Chapter~\ref{c5:ch}'s \code{sync}, synchronizing the \code{rx} line, which arrives from
    somebody else's crystal and is asynchronous in exactly the sense Chapter~\ref{c4:ch} means.
  \item Chapter~\ref{c7:ch}'s \code{timer}, generating the oversample tick.
  \item This chapter's state machine, framing the whole thing, and doing its own shifting and bit
    counting. You built that as a module in Chapter~\ref{c6:ch}; here it is four lines inside the
    machine that already has the counters it needs, which is the more usual way a receiver is
    written.
\end{itemize}
Write an entity named \code{uart_rx8}: a receiver for one byte of asynchronous serial data, in a
frame format modelled on a UART's, with one deliberate difference set out below.

\task{The line format}
The \code{rx} line carries frames, and between frames it sits idle:
\begin{itemize}
  \item \textbf{Idle}: the line is held high.
  \item \textbf{Start bit}: the line goes low for exactly one bit period. This is what marks the
    beginning of a frame; there is no separate clock line, which is what ``asynchronous serial''
    means.
  \item \textbf{Eight data bits}, most significant first, one bit period each. This is the
    deliberate difference: see the note below the figure.
  \item \textbf{Stop bit}: the line returns high for one bit period.
\end{itemize}

\bookfigure{uart_frame}{A frame on the \code{rx} line: idle high, a start bit, eight data bits most
significant first, and a stop bit, with the receiver sampling in the middle of every
bit.}{c8:fig:uartframe}

\noindent The frame in Figure~\ref{c8:fig:uartframe} carries \code{10110010}, which is the first
byte the testbench sends. The arrows are the only eight moments the receiver actually looks at the
line: one and a half bit periods after the falling edge, then once per bit period after that.

The receiver has to recover the bit timing from the falling edge of the start bit alone. That is
the entire problem this exercise solves, and it is why the timer is here.

\textbf{One deliberate deviation from real UART:} a real UART sends its data bits \textbf{least}
significant first, as \secref{c6:a3} notes. This exercise sends them most significant first:
\begin{itemize}
  \item the framing and the timing, which are what these two exercises are actually about, are
    identical either way.
  \item it is the direction \secref{c6:a3}'s SIPO idiom shifts, so the receiver's shift is one line,
    and the transmitter in Exercise~\ref{c8:ex:7} is its mirror.
  \item the provided testbenches transmit and expect MSB-first to match, so the pair is
    self-consistent.
  \item to speak to a real UART you would flip the shift direction to \secref{c6:a2}'s mirror-image
    form and change nothing else. Worth knowing before you wire this to a PC's serial port and
    wonder why every byte arrives bit-reversed.
\end{itemize}

\task{Oversampling: sixteen ticks per bit}
The receiver has to find the middle of a bit whose start it can only infer. The technique is
\textbf{oversampling}: run the \code{timer} at \textbf{sixteen times} the baud rate, so there are
sixteen ticks per bit period, and count them.
\begin{itemize}
  \item The falling edge of the start bit resets the tick count to \code{0}. The timer runs
    continuously, before, during and after a frame; it is the \emph{count} that is reset, never the
    timer.
  \item \textbf{Tick 8} is half a bit later, the middle of the start bit. Check the line is
    \emph{still low} there. If it has gone high, that was a glitch and not a start bit, and the
    receiver returns to idle. That re-check is what stops noise on an idle line from being read as
    data.
  \item From there, every \textbf{sixteenth} tick is the middle of the next bit: sample the line and
    shift it in. Eight of those are the data bits, and the ninth is the stop bit.
\end{itemize}
Sampling mid-bit rather than at the edges is what gives the receiver its timing margin. The
transmitter's clock is a different physical crystal from yours, so the two drift apart over a frame;
a receiver that sampled near the boundaries would read the wrong bit as soon as they did. Sixteen
ticks is the conventional choice.

\task{The \code{uart_rx8} entity}
\begin{booktable}{@{}p{11em}p{3.8em}p{3.6em}L@{}}
  \thead{Generic} & \thead{Type} & \thead{Default} & \thead{Description} \\
  \midrule
  \code{OVERSAMPLE_TICK_COUNT} & \code{natural} & \code{325} & The \code{timer} tick count for
    \textbf{one sixteenth} of a bit period. 9600 baud at \code{50 MHz}; part \textbf{b)} asks you to
    derive it. \\
\end{booktable}

\begin{booktable}{@{}p{5.2em}p{4.6em}p{10.6em}L@{}}
  \thead{Port} & \thead{Direction} & \thead{Type} & \thead{Description} \\
  \midrule
  \code{clock} & in & \code{std_logic} & System clock. \\
  \code{reset_n} & in & \code{std_logic} & Asynchronous, active-low; goes through \code{reset_sync},
    as always. \\
  \code{rx} & in & \code{std_logic} & The serial line, straight off a pin and therefore
    \textbf{asynchronous}. Synchronize it to \code{rx_s2} before anything else reads it. \\
  \code{data_out} & out & \code{std_logic_vector(7 downto 0)} & The received byte, valid in the
    cycle \code{byte_valid} is high. \\
  \code{byte_valid} & out & \code{std_logic} & A one-cycle pulse: a complete, well-framed byte is on
    \code{data_out}. \\
  \code{frame_err} & out & \code{std_logic} & A one-cycle pulse: the stop bit was not high, so the
    frame is broken and no byte is handed over. \\
\end{booktable}

\task{What to build: the receiver}
\begin{parts}
  \item Design the state machine \textbf{on paper first}, exactly as \secref{c8:a2} had you do.
    Four states:
    \begin{itemize}
      \item \code{STATE_IDLE}, waiting for the line to go low.
      \item \code{STATE_START}, checking at tick 8 that the line is still low.
      \item \code{STATE_DATA}, sampling one bit every sixteen ticks, eight times.
      \item \code{STATE_STOP}, sampling the stop bit sixteen ticks after the last data bit and
        deciding between \code{byte_valid} and \code{frame_err}.
    \end{itemize}
    Draw it before writing any VHDL, and on the diagram:
    \begin{itemize}
      \item Give every state its \textbf{self-loop}: the arrow for ``the tick I am waiting for has
        not arrived yet, stay put''. That is the arrow that goes missing, and a state without it
        falls through to wherever the \code{else} happens to point.
      \item Write beside each state \textbf{how many ticks it waits for}, so the tick accounting is
        on the paper rather than in your head.
      \item Mark where \code{ticks} is \textbf{cleared}. The timer itself never stops, so those
        clears are the only places the receiver decides what ``the start of a bit'' means.
      \item Label the arrow out of \code{STATE_IDLE} with an \textbf{edge}, not a level:
        \code{rx_s2} low \emph{and} high on the previous tick. That needs one bit of history, so
        note where you are going to keep it.
    \end{itemize}
  \item Work out the timing before you implement it:
    \begin{itemize}
      \item The DE0-CV clock is \code{50 MHz} and 9600 baud means 9600 bits per second. How many
        clock cycles is one bit period? How many is a sixteenth of one?
      \item \code{timer} pulses every \code{TICK_COUNT + 1} clock cycles (\secref{c7:a1}). What
        \code{OVERSAMPLE_TICK_COUNT} does that give you? The exact answer is not a whole number, so
        there are two candidates either side of it: work out the error each one leaves, and say
        which the default above is and why.
      \item That error is a fraction of a percent, and it accumulates across a frame. The stop bit
        is sampled 9.5 bit periods after the start edge: half a bit period to reach the middle of
        the start bit, then nine more bit periods. How far has your sample point drifted from the
        middle of the stop bit by then, measured in oversample ticks? How many ticks of margin did
        it have to begin with?
      \item Now the reason the answer is comfortable: you resynchronize on every start bit. What
        would have to be true of the frame length for that to stop being enough?
    \end{itemize}
  \item Implement \code{uart_rx8}. It instantiates three modules and does the rest itself:
    \begin{itemize}
      \item \code{reset_sync} for the reset. Everything else in the design takes its
        \code{reset_s2_n}.
      \item \textbf{\code{sync}}, the module you wrote in Chapter~\ref{c5:ch}'s
        Exercise~\ref{c5:ex:4}, on the \code{rx} pin, giving the \code{rx_s2} the state machine
        reads. Instantiate it at \code{SIZE = 1} and \code{PRESET = '1'}: an idle line is high,
        and a synchronizer coming out of reset reading \code{'0'} would look to this machine like a
        start bit. \code{rx} arrives from a pin driven by a transmitter with its own crystal, so it
        is asynchronous in exactly the sense Chapter~\ref{c4:ch} means, and every rule from that
        chapter applies unchanged. Note it takes a one-element \code{std_logic_vector} either
        side, so you declare a pair of those and pull \code{rx_s2} out of the output.
      \item \code{timer}, with \code{OVERSAMPLE_TICK_COUNT} passed straight through, \code{enable}
        tied to \code{'1'}, and its \code{timeout} as your only sense of elapsed time. Call it
        \code{baud_tick}. The timer free-runs: it is never stopped, never cleared, and no state ever
        touches it. What the machine re-phases instead is its own \code{ticks} counter.
      \item No shift register module. The machine already needs a bit counter, so it shifts the
        line straight into a \code{frame} signal itself, one line of VHDL per bit sampled.
    \end{itemize}
    \textbf{The state machine, state by state.} Two counters: \code{ticks} counting \code{0} to
    \code{15} within one bit, and \code{bit_count} counting the eight data bits. Everything below
    happens \textbf{only when \code{baud_tick = '1'}}; on every other clock edge the machine sits
    still. Since the timer free-runs, a tick is always on its way, in every state including idle.

    \begin{booktable}{@{}p{6.6em}>{\raggedright\arraybackslash}p{5.8em}L
      >{\raggedright\arraybackslash}p{7.6em}@{}}
      \thead{State} & \thead{Waits for} & \thead{Then} & \thead{To} \\
      \midrule
      \code{STATE_IDLE} & falling edge on \code{rx_s2} & clear \code{ticks} & \code{STATE_START} \\
      \code{STATE_START} & \code{ticks} = 8 & \code{rx_s2} still low? clear \code{ticks} &
        \code{STATE_DATA} \\
      & & \code{rx_s2} high? it was a glitch & \code{STATE_IDLE} \\
      \code{STATE_DATA} & \code{ticks} = 15 & shift \code{rx_s2} into \code{frame}, clear
        \code{ticks} & \code{STATE_DATA} until 8 bits, then \code{STATE_STOP} \\
      \code{STATE_STOP} & \code{ticks} = 15 & \code{rx_s2} high? \code{data_out <= frame}, pulse
        \code{byte_valid} & \code{STATE_IDLE} \\
      & & \code{rx_s2} low? pulse \code{frame_err}, drop the byte & \code{STATE_IDLE} \\
    \end{booktable}

    There is no ``done'' state. The stop-bit sample is where the frame is judged, so that is where
    the byte is handed over, and the machine goes straight back to idle.

    \textbf{Why idle waits for an edge rather than a level.} Keep one tick of history for
    \code{rx_s2}, and leave idle only on the transition from high to low. Testing the level instead,
    ``if \code{rx_s2} is low, start a frame'', looks equivalent and passes everything except the one
    case that matters.

    Consider a \textbf{break}: the line held low far longer than a frame, which is what a
    transmitter losing power looks like. A level test re-arms the instant the previous frame's stop
    bit is judged, so the receiver spends the break marching through back-to-back all-zero frames.
    Each ends on a low stop bit and is correctly rejected, and that part is fine. The problem is the
    frame in flight when the break \emph{ends}: its stop bit is sampled after the line has returned
    high, so it is a perfectly well-formed frame as far as the receiver can tell, and a byte made of
    nothing is handed over as real data. The edge test does not re-arm during the break at all,
    because after the first falling edge there is not another one until the line has recovered.

    That is the whole difference, it costs one flip-flop, and it is why a receiver that ``works'' on
    clean data can still hand you garbage the first time a cable is unplugged.

    \textbf{The shift.} Eight data bits arrive most significant first, so the SIPO idiom from
    \secref{c6:a3} does it with no index arithmetic at all:

\begin{vhdlcode}
frame <= frame(6 downto 0) & rx_s2;
\end{vhdlcode}

    \code{bit_count} then only counts \emph{how many} bits have arrived, never \emph{where} they go.
    Storing with \code{frame(bit_count) <= rx_s2} instead puts the first-arriving bit in bit 0,
    which reverses the byte: that is the least-significant-first order a real UART uses. Here the
    line is MSB first, so the shift is the one that matches.

    Four things that decide whether this works:
    \begin{itemize}
      \item \textbf{Everything is gated on \code{baud_tick}.} A state that acts on a plain clock edge
        runs sixteen times per tick and roughly five thousand times per bit.
      \item \textbf{Every state needs its ``not yet'' path.} If the tick you are waiting for has not
        arrived, or the count has not reached its target, the machine must \textbf{stay where it
        is}. Sending it anywhere else in the \code{else} is the single most common way this design
        fails, and it fails completely rather than subtly: the machine ends up ping-ponging between
        two states and never reaches the third.
      \item \textbf{Watch the off-by-one on the counter.} With \code{ticks} starting at \code{0} and
        incrementing on each tick, the \emph{n}th tick is \code{ticks = n - 1} at the moment you
        test it. Whether you compare before or after incrementing decides which tick you actually
        sample on, and being one or two ticks late is survivable while being eight is not.
      \item \textbf{The timer free-runs, so your phase comes from \code{ticks}, never from the
        timer.} The start edge arrives whenever the transmitter sends it, which is somewhere inside
        a tick period the receiver did not choose, so the first tick after that edge can be anywhere
        from a full tick to almost no time later. Clearing \code{ticks} on the way out of idle is
        what re-phases sampling onto the start bit. That leftover fraction of one tick is the price
        of never stopping the timer, and it is cheap: one sixteenth of a bit, against the half-bit
        of margin part \textbf{b)} asks you to work out.
    \end{itemize}
  \item Run the testbench. It needs four design files, subcomponents first, and none of the three
    subcomponents is shipped in the exercise directory: all three are modules you wrote,
    \code{reset_sync} in Chapter~\ref{c4:ch}, \code{sync} in Chapter~\ref{c5:ch} and \code{timer}
    in Chapter~\ref{c7:ch}.

\begin{shell}
cd lectures/L08/exercises/uart_rx8
cp ../../../L04/exercises/reset_sync/reset_sync.vhd .     # the three you wrote yourself
cp ../../../L05/exercises/sync/sync.vhd .
cp ../../../L07/exercises/timer/timer.vhd .
ghdl -a --std=93 sync.vhd reset_sync.vhd timer.vhd uart_rx8.vhd uart_rx8_tb.vhd
ghdl -e --std=93 uart_rx8_tb
ghdl -r --std=93 uart_rx8_tb --assert-level=error --stop-time=10ms
\end{shell}

    It acts as the transmitter, and checks:
    \begin{itemize}
      \item two clean frames, arriving as the right bytes, with no \code{frame_err}.
      \item an idle line, before, between and after them, producing nothing.
      \item a frame whose stop bit is \textbf{low}: \code{frame_err} must pulse and no byte may be
        handed over.
      \item the line held low far longer than a frame, which is a \textbf{break}: every frame the
        receiver thinks it sees there ends in a low stop bit, so none of them is a byte.
      \item a frame in which each data bit carries its value only in the \textbf{middle} of the bit
        period and the opposite level either side. A receiver sampling near tick 8 reads it
        perfectly; one sampling near a boundary reads the neighbouring bit.
      \item that \code{byte_valid} and \code{frame_err} are each never high for two cycles in a
        row.
    \end{itemize}
    \note[Tip] build it in stages and run the testbench after each, rather than writing all four
    states and debugging them at once:
    \begin{itemize}
      \item First get the machine to leave idle on a start bit and reach \code{STATE_DATA}. If it
        never gets there, the ``not yet'' path is what to look at.
      \item Then sample one data bit and check it is the right one before worrying about the other
        seven.
      \item Then the stop bit and the two strobes.
    \end{itemize}
  \item Now break it deliberately, and think about what the result does and doesn't prove:
    \begin{itemize}
      \item Remove the stop-bit state, handing the byte over as soon as the eighth data bit is in.
      \item Predict what should go wrong, then run it. Two checks catch this, and it is worth
        knowing which fires first: the frame with a low stop bit no longer raises \code{frame_err},
        and it hands over a byte that should have been dropped.
      \item Now the other half of the lesson. Passing is the weaker signal, so find a break the
        testbench does \textbf{not} catch. Here is one: \textbf{delete the mid-start-bit re-check},
        so \code{STATE_START} moves to \code{STATE_DATA} at tick 8 without confirming the line is
        still low. The testbench never puts a glitch on an idle line, so nothing notices, and on
        real hardware a single noise spike starts a phantom frame. Read \file{uart_rx8_tb.vhd} and
        say what you would add to catch it.
      \item Here is a second: \textbf{sample at tick 3 instead of tick 8.} Work out why the
        testbench still passes, and what you have given away. The answer is in part \textbf{b)}'s
        margin calculation.
      \item And a third, which is the most important of the three: \textbf{delete the \code{rx}
        synchronizer} and read the pin directly. Every check still passes, and it always will, no
        matter what is added to the testbench. Say why, in terms of what a simulator models and
        what it does not, then reread \secref{c4:a4}. This is the one class of bug in the book
        that testing cannot reach, which is exactly why Chapter~\ref{c4:ch} gives you a rule to
        follow rather than a symptom to look for.
      \item In a few sentences: what does ``passes its testbench'' actually promise, and what does
        it not?
    \end{itemize}
\end{parts}

\modulefigure{uart_rx8}

\selfcheck name your entity \code{uart_rx8}, with generic \code{OVERSAMPLE_TICK_COUNT}
(\code{natural}, default \code{325}), inputs \code{clock}, \code{reset_n}, \code{rx}, and outputs
\code{data_out} (\code{std_logic_vector(7 downto 0)}), \code{byte_valid} and \code{frame_err},
declared in that order. Its testbench is in \file{exercises/uart_rx8}, and nothing else is:
\textbf{copy in your own \code{reset_sync}, \code{sync} and \code{timer}}, as part \textbf{d)}
shows. A project directory holds everything it builds from, and assembling it is part of the
exercise. Chapter~\ref{c7:ch}'s \code{walking_led} is put together the same way.

\exercise{The serial transmitter}{Paper, VHDL}
\label{c8:ex:7}
The other half of the link: write \code{uart_tx8}, a transmitter that sends one byte in the same
format Exercise~\ref{c8:ex:6} receives.

A transmitter is the easier of the two, and the reason is worth understanding before you start:
\begin{itemize}
  \item The receiver had to \emph{recover} the bit timing from a start edge it didn't control, which
    is why it sampled mid-bit and why its timer ran at a sixteenth of a bit period.
  \item The transmitter \textbf{defines} the timing. Nothing has to be recovered, so the timer runs
    at a full bit period and the machine simply holds each bit on the line for one timeout.
\end{itemize}
It sends most significant bit first, matching Exercise~\ref{c8:ex:6}'s receiver rather than a real
UART, for the reason given there.

\task{The \code{uart_tx8} entity}
\begin{booktable}{@{}p{7.4em}p{3.8em}p{3.6em}L@{}}
  \thead{Generic} & \thead{Type} & \thead{Default} & \thead{Description} \\
  \midrule
  \code{BIT_TICK_COUNT} & \code{natural} & \code{5207} & The \code{timer} tick count for one
    \textbf{whole} bit period, 9600 baud at \code{50 MHz}. Sixteen of Exercise~\ref{c8:ex:6}'s
    oversample ticks, to within the rounding its part \textbf{b)} asks you to work out, so a
    \code{uart_tx8} and a \code{uart_rx8} built with the matching defaults speak at the same
    rate. \\
\end{booktable}

\begin{booktable}{@{}p{5.2em}p{4.6em}p{10.6em}L@{}}
  \thead{Port} & \thead{Direction} & \thead{Type} & \thead{Description} \\
  \midrule
  \code{clock} & in & \code{std_logic} & System clock. \\
  \code{reset_n} & in & \code{std_logic} & Asynchronous, active-low; through \code{reset_sync}, as
    always. \\
  \code{data_in} & in & \code{std_logic_vector(7 downto 0)} & The byte to send. \\
  \code{send} & in & \code{std_logic} & A request to transmit \code{data_in}. Ignored while the
    transmitter is busy. \\
  \code{tx} & out & \code{std_logic} & The serial line. \textbf{Idles high}, including during
    reset. \\
  \code{busy} & out & \code{std_logic} & High from the moment a frame starts until the stop bit is
    over. \\
\end{booktable}

\task{What to build: the transmitter}
\begin{parts}
  \item Draw the state machine first. Four states is the natural decomposition: idle, start bit,
    data bits, stop bit, and each one holds the line at a defined level for one bit period:
    \begin{itemize}
      \item Mark what drives \code{tx} in each state. Notice that \code{tx} never depends on
        \code{send} or any other primary input: it is the state, plus the shift register's current
        output. That makes this a Moore machine, like everything else in this book.
      \item Mark where \code{busy} comes from. It is one comparison against the state.
      \item Decide where the byte is captured. It has to be latched when the frame starts, not read
        continuously from \code{data_in}: the producer is free to change \code{data_in} the moment
        \code{busy} goes high, and a real one will.
    \end{itemize}
  \item Shift the byte out most significant bit first, in the machine itself, the mirror of what
    the receiver does:
    \begin{itemize}
      \item Latch \code{data_in} into a \code{frame} signal when the frame starts.
      \item Drive \code{tx} from the top bit of \code{frame} while in the data state.
      \item Shift \code{frame} up one place per bit period, and count the eight bits out.
    \end{itemize}
    You wrote this as a module in Chapter~\ref{c6:ch}'s Exercise~\ref{c6:ex:7}~b), and \code{piso8}
    is worth rereading for the idiom. Instantiating it here would work, but the machine already
    owns the bit counter that a PISO register needs, so the shift is one line inside the state it
    belongs to. That is how both halves of this link are built.
  \item Implement it, instantiating \code{reset_sync} for the reset and \code{timer} with
    \code{BIT_TICK_COUNT} passed straight through.

    Here the timer's \code{enable} \textbf{is} driven by the state, high only while a frame is in
    flight, and that is the one place this design deliberately differs from
    Exercise~\ref{c8:ex:6}. The receiver had to leave its timer free-running because it does not
    control when a frame starts: it can only re-phase its own tick count once the start edge has
    already arrived. The transmitter has the opposite problem. It \emph{decides} when the frame
    starts, so it can start the timer at that moment and get a first bit period exactly as long as
    every other one. Free-run it here instead and the start bit is short by however much of a tick
    period happened to be left when \code{send} arrived, which is a defect in the one thing a
    transmitter is responsible for.

    Two requirements the testbench is strict about:
    \begin{itemize}
      \item \textbf{\code{send} is ignored while busy.} Reloading mid-frame would restart the byte
        and corrupt what is already on the wire. A request that arrives during a frame is dropped,
        not queued.
      \item \textbf{\code{busy} stays high through the whole stop bit}, not just the data bits. The
        frame is not over until the line has been held high for that last bit period.
    \end{itemize}
  \item Run the testbench. It acts as the receiver: it waits for your start bit, samples the middle
    of each bit period, and checks the framing, the byte, the \code{busy} handshake, and that a
    \code{send} injected two bit periods into a frame is ignored.

\begin{shell}
cd lectures/L08/exercises/uart_tx8
cp ../../../L04/exercises/reset_sync/reset_sync.vhd .     # the two you wrote yourself
cp ../../../L07/exercises/timer/timer.vhd .
ghdl -a --std=93 reset_sync.vhd timer.vhd uart_tx8.vhd uart_tx8_tb.vhd
ghdl -e --std=93 uart_tx8_tb
ghdl -r --std=93 uart_tx8_tb --assert-level=error --stop-time=10ms
\end{shell}
  \item Answer, in prose:
    \begin{itemize}
      \item A transmitter that skips the stop bit entirely still leaves the line high afterwards,
        because idle and stop are the same level. What actually goes wrong, and when does it first
        show up?
        \begin{itemize}
          \item The testbench catches this one. Which of its checks does it, and why could no
            amount of watching \code{tx} alone have caught it?
        \end{itemize}
      \item Your receiver in Exercise~\ref{c8:ex:6} samples in the middle of each bit; your
        transmitter changes the line at the edges. If both ran off crystals differing by 1\%,
        roughly how many bits into a frame would the receiver's sampling point drift into the wrong
        bit? What does that tell you about why serial formats have a start bit per byte, rather than
        one at the beginning of a long message?
    \end{itemize}
\end{parts}

\modulefigure{uart_tx8}

\selfcheck name your entity \code{uart_tx8}, with generic \code{BIT_TICK_COUNT} (\code{natural},
default \code{5207}), inputs \code{clock}, \code{reset_n}, \code{data_in}
(\code{std_logic_vector(7 downto 0)}), \code{send}, and outputs \code{tx}, \code{busy}, declared in
that order. Its testbench is in \file{exercises/uart_tx8}, and nothing else is: copy in your own
\code{reset_sync} and \code{timer}, as part \textbf{d)} shows.

